\documentclass[11pt,a4paper]{article}

\usepackage[margin=1.6cm]{geometry}
\usepackage[T1]{fontenc}
\usepackage[utf8]{inputenc}
\usepackage{lmodern}
\usepackage{xcolor}
\usepackage{tabularx}
\usepackage{booktabs}
\usepackage{array}
\usepackage{amssymb}
\usepackage{enumitem}
\usepackage{hyperref}
\usepackage{fancyhdr}
\usepackage{titlesec}
\usepackage{tcolorbox}

\definecolor{pmoTeal}{HTML}{2F6F73}
\definecolor{pmoDark}{HTML}{182126}
\definecolor{pmoPink}{HTML}{D81E5B}
\definecolor{pmoAmber}{HTML}{D99A2B}
\definecolor{pmoLight}{HTML}{FDEDF4}

\hypersetup{
  colorlinks=true,
  urlcolor=pmoTeal,
  pdftitle={Sponsor Decision Forum Discussion Exercise},
  pdfsubject={Discussion exercise for From Plan to Control},
  pdfauthor={Mohamad Aoude}
}

\pagestyle{fancy}
\setlength{\headheight}{14pt}
\fancyhf{}
\lhead{From Plan to Control}
\rhead{Discussion Exercise}
\cfoot{\thepage}
\renewcommand{\headrulewidth}{0.4pt}

\titleformat{\section}{\Large\bfseries\color{pmoPink}}{}{0pt}{}
\titleformat{\subsection}{\large\bfseries\color{pmoDark}}{}{0pt}{}
\setlist[itemize]{leftmargin=1.4em,itemsep=2pt,topsep=4pt}
\setlist[enumerate]{leftmargin=1.7em,itemsep=3pt,topsep=4pt}
\setlength{\parindent}{0pt}
\setlength{\parskip}{5pt}

\newcolumntype{Y}{>{\raggedright\arraybackslash}X}
\newcommand{\blankline}{\rule{\linewidth}{0.35pt}}

\newtcolorbox{forumbox}{
  colback=pmoLight,
  colframe=pmoPink,
  boxrule=0.8pt,
  arc=2mm,
  left=3mm,right=3mm,top=2mm,bottom=2mm
}

\newtcolorbox{rulebox}{
  colback=yellow!8,
  colframe=pmoAmber,
  boxrule=0.8pt,
  arc=2mm,
  left=3mm,right=3mm,top=2mm,bottom=2mm
}

\begin{document}

\begin{center}
  {\small\bfseries\color{pmoPink} DISCUSSION ACTIVITY}\par
  \vspace{2mm}
  {\Huge\bfseries\color{pmoDark} Sponsor Decision Forum}\par
  \vspace{2mm}
  {\Large What Should the PMO Protect?}\par
  \vspace{3mm}
  {\small Duration: 20 minutes \quad | \quad Format: structured class debate \quad | \quad Evidence: position and final judgment}
\end{center}

\begin{forumbox}
\textbf{Discussion challenge.} The Smart Campus project is late, over budget, and facing a database-integration failure. There is no perfect response. Your task is to defend one recovery priority, question competing proposals, and revise or confirm your judgment after hearing evidence from others.
\end{forumbox}

\section{1. Learning Outcomes}

By the end of the forum, students will be able to:

\begin{itemize}
  \item state a clear project-recovery position;
  \item support claims using schedule, cost, scope, quality, and risk evidence;
  \item ask constructive questions and respond to counterarguments;
  \item compare the consequences of competing stakeholder priorities; and
  \item reach an individual final judgment after considering peer perspectives.
\end{itemize}

\section{2. Evidence Available to Everyone}

\begin{center}
\begin{tabularx}{\textwidth}{>{\bfseries}p{3.5cm} Y}
\toprule
Project evidence & Current situation \\
\midrule
Deadline & Pilot required at the end of week 12; current time is week 8. \\
Budget & Approved budget is USD 40,000; spending is already above the week 6 plan. \\
Progress & Actual progress was 38\% when planned progress was 50\%. \\
Quality & Twelve critical bugs were open; integration failure now blocks live data. \\
Essential pilot scope & Secure login, course schedules, and essential notifications. \\
Optional pilot scope & Room reservations, complaint workflow, analytics, and visual customization. \\
Team condition & Technical staff report overload and increasing rework. \\
Main uncertainty & The complete effort needed to stabilize integration is not yet known. \\
\bottomrule
\end{tabularx}
\end{center}

\section{3. Choose an Initial Position (2 minutes)}

Select one position before the discussion begins. You may change your position later, but you must explain why.

\begin{itemize}
  \item[$\square$] \textbf{Position A -- Protect the deadline:} freeze optional scope and deliver a smaller reliable pilot in week 12.
  \item[$\square$] \textbf{Position B -- Protect the complete scope:} request a later deadline so all promised functions remain in the first pilot.
  \item[$\square$] \textbf{Position C -- Protect scope and deadline through investment:} request additional budget for focused integration expertise.
  \item[$\square$] \textbf{Position D -- Protect the approved budget:} retain current resources and accept a reduced scope or later completion based on a new forecast.
\end{itemize}

\textbf{My initial position:} \blankline

\textbf{My strongest supporting evidence:}

\vspace{8mm}\blankline

\textbf{The stakeholder most helped by this position:} \blankline

\textbf{The stakeholder or objective placed most at risk:} \blankline

\section{4. Discussion Rules}

\begin{rulebox}
\begin{enumerate}
  \item Critique ideas and evidence, never the person presenting them.
  \item Use the structure \textbf{claim $\rightarrow$ evidence $\rightarrow$ reasoning}.
  \item Do not repeat a previous point unless you add new evidence or a different implication.
  \item Before disagreeing, summarize the previous speaker's point fairly.
  \item The goal is not to win. The goal is to improve the quality of the project decision.
\end{enumerate}
\end{rulebox}

\section{5. Forum Sequence (15 minutes)}

\subsection*{Round 1 -- Opening positions (4 minutes)}

One representative for each position gives a maximum 45-second statement:

\begin{forumbox}
``The PMO should \underline{\hspace{4cm}} because the evidence shows \underline{\hspace{4cm}}. This protects \underline{\hspace{3cm}}, although it risks \underline{\hspace{3cm}}.''
\end{forumbox}

Record one strong point from a position different from yours:

\vspace{8mm}\blankline

\subsection*{Round 2 -- Evidence questions (4 minutes)}

Ask questions that expose assumptions or missing evidence. Useful prompts include:

\begin{itemize}
  \item What evidence supports your forecast?
  \item Which stakeholder accepts the consequence of this choice?
  \item What happens if integration takes longer than expected?
  \item How does your proposal protect quality and security?
  \item What exact sponsor approval does your proposal require?
\end{itemize}

\textbf{Question I asked or wanted to ask:}

\vspace{7mm}\blankline

\textbf{Most useful answer heard:}

\vspace{7mm}\blankline

\subsection*{Round 3 -- Counterargument and response (4 minutes)}

Complete both statements:

\textbf{A reasonable objection to my position is:}

\vspace{7mm}\blankline

\textbf{My evidence-based response is:}

\vspace{7mm}\blankline

\subsection*{Round 4 -- Final vote (3 minutes)}

Vote again after hearing the discussion. The final vote is individual; consensus is not required.

\begin{itemize}
  \item[$\square$] Position A
  \item[$\square$] Position B
  \item[$\square$] Position C
  \item[$\square$] Position D
\end{itemize}

\section{6. Final Judgment}

\textbf{My final position:} \blankline

\textbf{Did my position change? Why or why not?}

\vspace{9mm}\blankline

\textbf{The most persuasive evidence or argument was:}

\vspace{9mm}\blankline

\textbf{The sponsor decision I would request is:}

\vspace{9mm}\blankline

\section{7. Discussion Assessment}

\begin{center}
\begin{tabularx}{\textwidth}{>{\bfseries}p{4cm} Y >{\centering\arraybackslash}p{2cm}}
\toprule
Criterion & Observable evidence & Points \\
\midrule
Clear position & States a specific, relevant recovery priority. & 3 \\
Use of evidence & Uses project data rather than unsupported opinion. & 3 \\
Reasoning & Explains consequences and recognizes trade-offs. & 3 \\
Dialogic participation & Listens, questions, and responds constructively. & 3 \\
Final judgment & Revises or confirms the position using ideas from the discussion. & 3 \\
\midrule
Total & & 15 \\
\bottomrule
\end{tabularx}
\end{center}

\section{8. Instructor Facilitation Notes}

\begin{itemize}
  \item Display or distribute the shared evidence before students choose positions.
  \item Place students by initial position, then alternate speakers across positions.
  \item Ask quiet students to contribute a question before allowing a second contribution from frequent speakers.
  \item Do not announce one universally correct answer. Evaluate the evidence, assumptions, trade-offs, and requested sponsor decision.
  \item Close by distinguishing \textbf{discussion} from \textbf{collaboration}: this forum improves individual judgment through exchanged perspectives; the collaboration exercise requires a team to produce one shared recovery briefing.
\end{itemize}

\section{9. Responsible AI Boundary}

Students must form an initial position before using AI. If permitted, AI may generate a counterargument that students critically evaluate. It may not vote, invent project evidence, or replace peer discussion. No personal or confidential information may be uploaded.

\vfill
\begin{center}
\small\color{pmoTeal}
Course portal: \href{https://maoude.github.io/teaching-ai-pmo/}{maoude.github.io/teaching-ai-pmo/}
\end{center}

\end{document}
